\documentclass[11pt]{article}
\usepackage[margin=1.1in]{geometry}
\usepackage{amsmath,amssymb,amsthm,mathtools,booktabs,array,tabularx,longtable}
\usepackage[dvipsnames]{xcolor}
\usepackage[colorlinks=true,linkcolor=MidnightBlue,citecolor=MidnightBlue,urlcolor=MidnightBlue]{hyperref}
\usepackage{fancyhdr}
\pagestyle{fancy}\fancyhf{}\renewcommand{\headrulewidth}{0pt}
\fancyfoot[R]{\thepage}
\fancyfoot[L]{\scriptsize\color{gray}Generated by the Transport Lab from approved parts of its store; not a deposited paper. Papers of record: \href{https://zenodo.org/communities/transport-papers}{Zenodo, Transport Papers}}
\fancypagestyle{plain}{\fancyhf{}\fancyfoot[R]{\thepage}\fancyfoot[L]{\scriptsize\color{gray}Generated by the Transport Lab from approved parts of its store; not a deposited paper. Papers of record: \href{https://zenodo.org/communities/transport-papers}{Zenodo, Transport Papers}}\renewcommand{\headrulewidth}{0pt}}
\newtheorem{theorem}{Theorem}[section]
\newtheorem{proposition}[theorem]{Proposition}
\newtheorem{lemma}[theorem]{Lemma}
\newtheorem{corollary}[theorem]{Corollary}
\newtheorem{conjecture}[theorem]{Conjecture}
\newtheorem{refutation}[theorem]{Refutation}
\newtheorem{prediction}[theorem]{Prediction}
\newtheorem{principle}[theorem]{Principle}
\newtheorem{observation}[theorem]{Observation}
\newtheorem{axiom}[theorem]{Axiom}
\theoremstyle{definition}
\newtheorem{definition}[theorem]{Definition}
\newtheorem{fact}[theorem]{Fact}
\newtheorem{assembled}[theorem]{Assembly}
\newtheorem{postulate}[theorem]{Postulate}
\newtheorem{premise}[theorem]{Premise}
\newtheorem{hypothesis}[theorem]{Hypothesis}
\newtheorem{remark}[theorem]{Remark}
\newtheorem{example}[theorem]{Example}
\theoremstyle{plain}
\setlength{\marginparwidth}{0.8in}
\newcommand{\changedmark}[1]{\marginpar{\raggedright\scriptsize\color{BrickRed}#1}}
\providecommand{\Aa}{}\renewcommand{\Aa}{\mathcal{A}}
\providecommand{\Ad}{}\renewcommand{\Ad}{\widetilde{\mathrm{Ad}}}
\providecommand{\Aev}{}\renewcommand{\Aev}{\mathcal{A}^{+}}
\providecommand{\Aod}{}\renewcommand{\Aod}{\mathcal{A}^{-}}
\providecommand{\B}{}\renewcommand{\B}{\mathrm{B}}
\providecommand{\C}{}\renewcommand{\C}{\mathbb{C}}
\providecommand{\CAR}{}\renewcommand{\CAR}{\mathrm{CAR}}
\providecommand{\CS}{}\renewcommand{\CS}{{\mathbb{S}^{+}}}
\providecommand{\CSm}{}\renewcommand{\CSm}{{\mathbb{S}^{-}}}
\providecommand{\Cl}{}\renewcommand{\Cl}{\mathrm{Cl}}
\providecommand{\Con}{}\renewcommand{\Con}{\mathcal{C}}
\providecommand{\Dl}{}\renewcommand{\Dl}{\Delta}
\providecommand{\EN}{}\renewcommand{\EN}{E_{N}}
\providecommand{\End}{}\renewcommand{\End}{\mathrm{End}}
\providecommand{\F}{}\renewcommand{\F}{\mathbb{F}}
\providecommand{\G}{}\renewcommand{\G}{\Gamma}
\providecommand{\HH}{}\renewcommand{\HH}{\mathbb{H}}
\providecommand{\IS}{}\renewcommand{\IS}{I\sigma_3}
\providecommand{\J}{}\renewcommand{\J}{\mathbf{J}}
\providecommand{\LH}{}\renewcommand{\LH}{L_{H}}
\providecommand{\Mat}{}\renewcommand{\Mat}{\mathrm{M}}
\providecommand{\Mo}{}\renewcommand{\Mo}{M^{\circ}}
\providecommand{\Mp}{}\renewcommand{\Mp}{M^{+}}
\providecommand{\Msun}{}\renewcommand{\Msun}{M_{\odot}}
\providecommand{\Mt}{}\renewcommand{\Mt}{\widetilde{M}}
\providecommand{\Ncone}{}\renewcommand{\Ncone}{\mathcal{N}}
\providecommand{\Neff}{}\renewcommand{\Neff}{N_{\mathrm{eff}}}
\providecommand{\PCS}{}\renewcommand{\PCS}{\mathbb{P}{\mathbb{S}^{+}}}
\providecommand{\PT}{}\renewcommand{\PT}{\mathbb{PT}}
\providecommand{\Pin}{}\renewcommand{\Pin}{\mathrm{Pin}}
\providecommand{\Q}{}\renewcommand{\Q}{\mathbb{Q}}
\providecommand{\R}{}\renewcommand{\R}{\mathbb{R}}
\providecommand{\Ric}{}\renewcommand{\Ric}{\mathrm{Ric}}
\providecommand{\SBH}{}\renewcommand{\SBH}{S_{\mathrm{BH}}}
\providecommand{\Sc}{}\renewcommand{\Sc}{\mathbb{S}}
\providecommand{\Sch}{}\renewcommand{\Sch}{\mathrm{P}}
\providecommand{\Sig}{}\renewcommand{\Sig}{\Sigma}
\providecommand{\Som}{}\renewcommand{\Som}{S_\omega}
\providecommand{\Sp}{}\renewcommand{\Sp}{\mathbb{S}}
\providecommand{\Sph}{}\renewcommand{\Sph}{\mathrm{S}}
\providecommand{\Spi}{}\renewcommand{\Spi}{S_\pi}
\providecommand{\Spin}{}\renewcommand{\Spin}{\mathrm{Spin}}
\providecommand{\Sth}{}\renewcommand{\Sth}{S_{\mathrm{th}}}
\providecommand{\T}{}\renewcommand{\T}{\mathbb{T}}
\providecommand{\TH}{}\renewcommand{\TH}{T_{\mathrm{H}}}
\providecommand{\Tr}{}\renewcommand{\Tr}{\operatorname{Tr}}
\providecommand{\UHF}{}\renewcommand{\UHF}{\mathrm{UHF}}
\providecommand{\Z}{}\renewcommand{\Z}{\mathbb{Z}}
\providecommand{\Zc}{}\renewcommand{\Zc}{\mathrm{Z}}
\providecommand{\cmg}{}\renewcommand{\cmg}{\mathrm{cm^{2}\,g^{-1}}}
\providecommand{\dS}{}\renewcommand{\dS}{\mathrm{dS}}
\providecommand{\dsc}{}\renewcommand{\dsc}{\mathrm{dsc}}
\providecommand{\dtPG}{}\renewcommand{\dtPG}{dt}
\providecommand{\dtS}{}\renewcommand{\dtS}{dt_{\mathrm{S}}}
\providecommand{\ellv}{}\renewcommand{\ellv}{\ell}
\providecommand{\er}{}\renewcommand{\er}{e_r}
\providecommand{\et}{}\renewcommand{\et}{e_t}
\providecommand{\ev}{}\renewcommand{\ev}{\mathrm{ev}}
\providecommand{\gl}{}\renewcommand{\gl}{\mathfrak{gl}}
\providecommand{\hatV}{}\renewcommand{\hatV}{\widehat{V}}
\providecommand{\hb}{}\renewcommand{\hb}{\bar{h}}
\providecommand{\hf}{}\renewcommand{\hf}{\hat{f}}
\providecommand{\id}{}\renewcommand{\id}{\mathrm{id}}
\providecommand{\kpc}{}\renewcommand{\kpc}{\mathrm{kpc}}
\providecommand{\lP}{}\renewcommand{\lP}{\ell_{P}}
\providecommand{\mD}{}\renewcommand{\mD}{\mathbf{m}_{D}}
\providecommand{\mb}{}\renewcommand{\mb}{\bar m}
\providecommand{\mnu}{}\renewcommand{\mnu}{\Sigma m_{\nu}}
\providecommand{\nInf}{}\renewcommand{\nInf}{n_\infty}
\providecommand{\nO}{}\renewcommand{\nO}{n_0}
\providecommand{\ninf}{}\renewcommand{\ninf}{n_{\infty}}
\providecommand{\nz}{}\renewcommand{\nz}{n_{0}}
\providecommand{\proj}{}\renewcommand{\proj}{\operatorname{proj}}
\providecommand{\pv}{}\renewcommand{\pv}{\pi}
\providecommand{\rad}{}\renewcommand{\rad}{\mathrm{rad}}
\providecommand{\sgn}{}\renewcommand{\sgn}{\operatorname{sgn}}
\providecommand{\shA}{}\renewcommand{\shA}{\mathcal{A}}
\providecommand{\shC}{}\renewcommand{\shC}{\mathcal{C}}
\providecommand{\shE}{}\renewcommand{\shE}{\mathcal{E}}
\providecommand{\shJ}{}\renewcommand{\shJ}{\mathcal{J}}
\providecommand{\shK}{}\renewcommand{\shK}{\mathcal{K}}
\providecommand{\shL}{}\renewcommand{\shL}{L}
\providecommand{\shW}{}\renewcommand{\shW}{\mathcal{W}}
\providecommand{\sll}{}\renewcommand{\sll}{\mathfrak{sl}}
\providecommand{\so}{}\renewcommand{\so}{\mathfrak{so}}
\providecommand{\su}{}\renewcommand{\su}{\mathfrak{su}}
\providecommand{\supp}{}\renewcommand{\supp}{\operatorname{supp}}
\providecommand{\tr}{}\renewcommand{\tr}{\operatorname{tr}}
\providecommand{\uu}{}\renewcommand{\uu}{\mathfrak{u}}
\providecommand{\br}{}\renewcommand{\br}[1]{\langle #1\rangle}
\providecommand{\bra}{}\renewcommand{\bra}[1]{\langle #1|}
\providecommand{\file}{}\renewcommand{\file}[1]{\texttt{#1}}
\providecommand{\hp}{}\renewcommand{\hp}[2]{( #1 \,|\, #2 )}
\providecommand{\ip}{}\renewcommand{\ip}[2]{\langle #1, #2\rangle}
\providecommand{\ket}{}\renewcommand{\ket}[1]{|#1\rangle}
\providecommand{\rel}{}\renewcommand{\rel}[2]{D\bigl(#1 \,\|\, #2\bigr)}
\providecommand{\rev}{}\renewcommand{\rev}[1]{\widetilde{#1}}
\providecommand{\slashed}{}\renewcommand{\slashed}[1]{{#1}\mkern-10mu/}
\title{\bfseries Two Rungs Up: The Composite as a Spinor of the Transported Algebra, and What That Would Decide: A companion note to Transport Papers T12 (Hypothesis (S), successor problem (6))}
\author{Leon Butler\\ \normalsize\href{mailto:lbutler.math@gmail.com}{\texttt{lbutler.math@gmail.com}}\\ \normalsize ORCID: \href{https://orcid.org/0009-0006-4449-5617}{0009-0006-4449-5617}}
\date{Generated by the Transport Lab, 27 October 2026}
\begin{document}
\maketitle
\begin{abstract}
\normalsize
The Standard-Model paper of the series needs one hypothesis for the lepton doublet, the singlet listing and its unifying group: that the composite's state space is the five-mode Fock space $\Lambda^{\bullet}(W \oplus \C^{2})$, three pair modes and two slot modes. This note asks whether the tower derives it. The exact part: the spinor module of the $n$-rung transport $\Cl(2,4) \to \Cl(2{+}n, 4{+}n)$ is $\Lambda^{\bullet}(W \oplus \C^{n})$, so the hypothesis says that a two-twistor composite lives two rungs up; the rung-and-pair mixings are $U(n) \times U(3)$ acting as the hypothesis prescribes; the ladder group of the Standard-Model paper is this $U(n)$ on the one-constituent sector; the compact $\mathrm{Spin}(10)$ is the real form fixed by the tower's adjoint structure once its frames are fixed; and the mass bracket of the mass paper is the frozen point's grip, at $\Lambda = 0$, on the slot-antisymmetric bivector. Given these and the identification of slots with rungs, the rung and slot-statistics hypotheses are one hypothesis in two languages. The part that does not close: no theorem of the series makes each twistor a rung. One completion transports every null momentum at once; the null bivectors of any number of twistors invite one further rung, together, and that rung's $U(2)$ mixes chirality; and the module holds the composite's tuple as one vector, not two twistors. So the hypothesis is sharpened to one structural sentence with a named missing mechanism and, failing that, registered at Tier C with priced kill conditions. Every identity is verified exactly by a deposited script.
\end{abstract}
\medskip
\noindent\textbf{MSC 2020.} Primary 81R40, 15A66; Secondary 81V22, 20G45.\\
\textbf{Keywords.} transport principle; Clifford tower; spinor module; Standard Model; $\mathrm{Spin}(10)$; twistor; Fock space.
\newpage
\noindent \textbf{Authorship and disclosure. }The founding idea of the series is the author's: that division by zero in a Clifford algebra should be read not as a prohibition but as a transport into a higher-dimensional algebra, and that the transported division need not be accompanied by an inverse operation. The direction of investigation and the critical review of the results are likewise the author's. The mathematical development and the drafting of the text were carried out in extended collaboration with Claude Fable 5 (Anthropic), an AI system.

\noindent\fbox{\parbox{\dimexpr\linewidth-2\fboxsep-2\fboxrule}{\textbf{This is not a deposited paper.} It was generated by the Transport Lab on 27 October 2026 and assembled from approved parts of its store. It follows results across several papers of the series; each is stated with its proof and credited to the paper that first proved it. The official versions of those papers are on Zenodo, in the Transport Papers community, \url{https://zenodo.org/communities/transport-papers}.}}
\newpage

\section{Setting, and what is imported}

\phantomsection\label{c:note.rungs.setting-setting-what-imported}
Notation is T12's \cite{ButlerT12}: the pair space $W$, three-dimensional and maximal isotropic in the complexified ambient space $V_{\C}$ of $\Cl(2,4)$; the Dirac module $\Sp = \Lambda^{\bullet}W$ with chirality halves $\T = \Lambda^{\mathrm{even}}W$ (twistors) and $\T' = \Lambda^{\mathrm{odd}}W$; the automorphism group $S(U(2) \times U(3))$ with its $U(1)$ factor $(\lambda^{3}a, \lambda^{-2}b)$; abelian weight $q$ and hypercharge $Y = q/6$. Hypothesis (S) of T12 is: the composite's state space is $\Lambda^{\bullet}(\C^{2} \oplus W)$, with $(A, B)$ acting as $\Lambda^{\bullet}(A \oplus B)$ and chirality the parity of the total mode number. Eight imports.

\begin{assembled}[Imported: the tower, one qubit per rung]\label{c:rungs.imported-tower-one-qubit-per-rung}
The transport principle completes $\Cl(p,q)$ to $\Cl(p{+}1,q{+}1)$ by adjoining one hyperbolic pair, and the completed algebra carries a null cone of its own, so the completion iterates: the $k$-th rung is $T_{k} = \Cl(p{+}k, q{+}k)$, with adjoined null pairs $v_{j}, w_{j}$ satisfying, as vectors, $v_{i}v_{j} + v_{j}v_{i} = 0$, $w_{i}w_{j} + w_{j}w_{i} = 0$ and $v_{i}w_{j} + w_{j}v_{i} = \delta_{ij}$. Inside $T_{k}$ the subalgebra $B_{k}$ generated by the $2k$ plane generators is $\Cl(k,k) \cong \mathrm{M}_{2^{k}}(\R)$, the twisted ground generators generate a commuting copy $A_{k} \cong \Cl(p,q)$, and multiplication $A_{k} \otimes B_{k} \to T_{k}$ is an isomorphism, $T_{k} \cong \Cl(p,q) \otimes \mathrm{M}_{2^{k}}(\R)$. With $W_{k} = \operatorname{span}\{v_{1}, \ldots, v_{k}\}$ the Chevalley action makes $\Lambda^{\bullet}W_{k}$, of dimension $2^{k}$, an irreducible $B_{k}$-module, the vacuum $1$ annihilated by every $w_{j}$ and cyclic for the $v_{j}$, and each rung doubles it: $\Lambda^{\bullet}W_{k+1} = \Lambda^{\bullet}W_{k} \otimes \Lambda^{\bullet}(\R v_{k+1})$, one qubit per rung. The Witt index rises by exactly one per rung, and the Witt index of a real quadratic space is the maximal number of anticommuting mode pairs realisable by vectors. On each complexified floor there is a unique conjugate-linear antiautomorphism $*$ fixing the positive frame generators and negating the negative ones, under which the ladder pairs are mutually adjoint, $v_{j}^{*} = w_{j}$. \emph{Provenance:} \cite[Def.~2.2, Thm.~3.2, Prop.~4.1, Thms.~4.2, 4.3, Lem.~5.4]{ButlerT3}.
\end{assembled}
\begin{proof}
The transport principle completes $\Cl(p,q)$ to $\Cl(p{+}1,q{+}1)$ by adjoining one hyperbolic pair, and the completed algebra carries a null cone of its own, so the completion iterates: the $k$-th rung is $T_{k} = \Cl(p{+}k, q{+}k)$, with adjoined null pairs $v_{j}, w_{j}$ satisfying, as vectors, $v_{i}v_{j} + v_{j}v_{i} = 0$, $w_{i}w_{j} + w_{j}w_{i} = 0$ and $v_{i}w_{j} + w_{j}v_{i} = \delta_{ij}$. Inside $T_{k}$ the subalgebra $B_{k}$ generated by the $2k$ plane generators is $\Cl(k,k) \cong \mathrm{M}_{2^{k}}(\R)$, the twisted ground generators generate a commuting copy $A_{k} \cong \Cl(p,q)$, and multiplication $A_{k} \otimes B_{k} \to T_{k}$ is an isomorphism, $T_{k} \cong \Cl(p,q) \otimes \mathrm{M}_{2^{k}}(\R)$. With $W_{k} = \operatorname{span}\{v_{1}, \ldots, v_{k}\}$ the Chevalley action makes $\Lambda^{\bullet}W_{k}$, of dimension $2^{k}$, an irreducible $B_{k}$-module, the vacuum $1$ annihilated by every $w_{j}$ and cyclic for the $v_{j}$, and each rung doubles it: $\Lambda^{\bullet}W_{k+1} = \Lambda^{\bullet}W_{k} \otimes \Lambda^{\bullet}(\R v_{k+1})$, one qubit per rung. The Witt index rises by exactly one per rung, and the Witt index of a real quadratic space is the maximal number of anticommuting mode pairs realisable by vectors. On each complexified floor there is a unique conjugate-linear antiautomorphism $*$ fixing the positive frame generators and negating the negative ones, under which the ladder pairs are mutually adjoint, $v_{j}^{*} = w_{j}$. \emph{Provenance:} \cite[Def.~2.2, Thm.~3.2, Prop.~4.1, Thms.~4.2, 4.3, Lem.~5.4]{ButlerT3} (\ref{c:tower.structure-tower-c-norms}, \ref{c:tower.witt-index-quantum-capacity}, \ref{c:tower.commuting-copies}, \ref{c:tower.fock-module-one-qubit-per-rung}, \ref{c:tower.ladder-pairs-satisfy-car-vectors}).

\end{proof}

\begin{assembled}[Imported: totality and minimality of the completion]\label{c:rungs.imported-totality-minimality-completion}
For a nondegenerate $(\mathbb{V}, Q)$ of mixed signature and dimension $p + q \geq 3$, no extension of $\Cl(\mathbb{V})$ carries a $\mathrm{Spin}^{0}$-equivariant element-level inverse of the null vectors; the minimal null-lift completion is the adjunction of a single hyperbolic plane, $\Cl(p,q) \to \Cl(p{+}1,q{+}1)$, unique up to isometry, in every signature; and on that completion's quadric the transported reciprocal $\iota$ is total, continuous and involutive, agreeing with the geometric reciprocal on invertible vectors and sending every null vector $x_{0}$ to the ideal point $[x_{0} + \tfrac12 n_{\infty}]$, as the limit along every invertible path. \emph{Provenance:} \cite[Thms.~4.2, 4.10, 4.13]{ButlerT1}.
\end{assembled}
\begin{proof}
For a nondegenerate $(\mathbb{V}, Q)$ of mixed signature and dimension $p + q \geq 3$, no extension of $\Cl(\mathbb{V})$ carries a $\mathrm{Spin}^{0}$-equivariant element-level inverse of the null vectors; the minimal null-lift completion is the adjunction of a single hyperbolic plane, $\Cl(p,q) \to \Cl(p{+}1,q{+}1)$, unique up to isometry, in every signature; and on that completion's quadric the transported reciprocal $\iota$ is total, continuous and involutive, agreeing with the geometric reciprocal on invertible vectors and sending every null vector $x_{0}$ to the ideal point $[x_{0} + \tfrac12 n_{\infty}]$, as the limit along every invertible path. \emph{Provenance:} \cite[Thms.~4.2, 4.10, 4.13]{ButlerT1} (\ref{c:transport.no-go-level-elements}, \ref{c:transport.minimality}, \ref{c:transport.totality}).

\end{proof}

\begin{assembled}[Imported: the spinor identification, the Witt data, and the line at infinity]\label{c:rungs.imported-spinor-identification-witt-data-line}
The complexified $\Cl(2,4)$ acts irreducibly on $\Sp = \Lambda^{\bullet}W$. There is a canonical isomorphism $V_{\C} \cong \Lambda^{2}\CS$ of $\mathrm{Spin}$-modules under which the quadratic form is $P \mapsto \varepsilon(P \wedge P)$ for a fixed volume functional $\varepsilon$ on $\Lambda^{4}\CS$, and a vector is null exactly when its bivector is decomposable. In the incidence chart $\CS = \Som \oplus \Spi$ of the conformal spinor space (twistor space), a conformal spinor being $Z = (\omega, \pi)$, the lines of $\PCS = \mathbb{P}(\CS)$ skew to the line at infinity $L_{\infty} = \mathbb{P}(\Som)$ are the points of the chart, and the Witt pair $n_{0}, n_{\infty}$ of $V_{\C}$ consists of the bivectors of the two summands, $\Som \wedge \Spi$ being their common orthogonal complement. The real form of the ambient spin group is the conformal $SU(2,2) \cong \mathrm{Spin}(2,4)$, and a line of $\PCS$ is a point of the real conformal compactification exactly when it is null for the Hermitian form $\Sigma$. \emph{Provenance:} \cite[Rem.~3.2, Thms.~3.1, 3.4, 4.2, Props.~4.1, 4.5]{ButlerT9}.
\end{assembled}
\begin{proof}
The complexified $\Cl(2,4)$ acts irreducibly on $\Sp = \Lambda^{\bullet}W$. There is a canonical isomorphism $V_{\C} \cong \Lambda^{2}\CS$ of $\mathrm{Spin}$-modules under which the quadratic form is $P \mapsto \varepsilon(P \wedge P)$ for a fixed volume functional $\varepsilon$ on $\Lambda^{4}\CS$, and a vector is null exactly when its bivector is decomposable. In the incidence chart $\CS = \Som \oplus \Spi$ of the conformal spinor space (twistor space), a conformal spinor being $Z = (\omega, \pi)$, the lines of $\PCS = \mathbb{P}(\CS)$ skew to the line at infinity $L_{\infty} = \mathbb{P}(\Som)$ are the points of the chart, and the Witt pair $n_{0}, n_{\infty}$ of $V_{\C}$ consists of the bivectors of the two summands, $\Som \wedge \Spi$ being their common orthogonal complement. The real form of the ambient spin group is the conformal $SU(2,2) \cong \mathrm{Spin}(2,4)$, and a line of $\PCS$ is a point of the real conformal compactification exactly when it is null for the Hermitian form $\Sigma$. \emph{Provenance:} \cite[Rem.~3.2, Thms.~3.1, 3.4, 4.2, Props.~4.1, 4.5]{ButlerT9} (\ref{c:correspondence.reality-conjugation-twistor-form-factor}, \ref{c:correspondence.fock-model-complexified-ascended-algebra}, \ref{c:correspondence.skew-lines-splittings-hyperbolic-planes-one}, \ref{c:correspondence.spinor-identification}, \ref{c:correspondence.incidence-relation-chart}).

\end{proof}

\begin{assembled}[Imported: mass as symplectic area]\label{c:rungs.imported-mass-symplectic-area}
On the real slice a single conformal spinor (twistor) carries null momentum $p_{AA'} = \pi_{A}\bar\pi_{A'}$, and a system of two conformal spinors with $p = \pi^{1}\bar\pi^{1} + \pi^{2}\bar\pi^{2}$ has $m^{2} = p^{2} = 2|\br{\pi^{1}\pi^{2}}|^{2}$, with $\br{\pi^{1}\pi^{2}} = \varepsilon^{AB}\pi^{1}_{A}\pi^{2}_{B}$. In the flat chart of the conformal spinor space $\CS$ the frozen point is $I_{\Lambda} = n_{\infty} + (\Lambda/6)\,n_{0}$, so at $\Lambda = 0$ it is $n_{\infty}$. \emph{Provenance:} \cite[Prop.~4.1, Lem.~5.1]{ButlerT11}; \cite[Fact~2.5]{ButlerT12}.
\end{assembled}
\begin{proof}
On the real slice a single conformal spinor carries null momentum $p_{AA'} = \pi_{A}\bar\pi_{A'}$, and a system of two conformal spinors with $p = \pi^{1}\bar\pi^{1} + \pi^{2}\bar\pi^{2}$ has $m^{2} = p^{2} = 2|\br{\pi^{1}\pi^{2}}|^{2}$, with $\br{\pi^{1}\pi^{2}} = \varepsilon^{AB}\pi^{1}_{A}\pi^{2}_{B}$. In the flat chart of $\CS$ the frozen point is $I_{\Lambda} = n_{\infty} + (\Lambda/6)\,n_{0}$, so at $\Lambda = 0$ it is $n_{\infty}$. \emph{Provenance:} \cite[Prop.~4.1, Lem.~5.1]{ButlerT11} (\ref{c:mass.de-sitter-scale-frozen-point}, \ref{c:mass.mass-symplectic-area-two-twistor-formula}); \cite[Fact~2.5]{ButlerT12} (\ref{c:standard-model.imported-frozen-point-dictionary}).

\end{proof}

\begin{assembled}[Imported: the ladder, the automorphism group, and the no-lepton clause]\label{c:rungs.imported-ladder-automorphism-group-no-lepton}
The mixings of a system of $n$ conformal spinors (an $n$-twistor system) preserving momentum and helicity for every configuration are exactly $U(n)$, acting on the $n$-tuple $(Z^{1}, \ldots, Z^{n}) \in \C^{n} \otimes \CS$, $\CS$ the conformal spinor space, as $U \otimes 1$; the automorphism group of the hierarchy $\C^{2} \otimes \C^{3}$ is $(U(2) \times U(3))/K \cong S(U(2) \times U(3))$, and that of $\C^{3} \otimes \C^{3}$ has Lie algebra $\su(3) \oplus \su(3) \oplus \uu(1)$, dimension seventeen; the tensor ledger $\Lambda^{\bullet}(\C^{2} \otimes \C^{3})$ contains no $(2,1)$. \emph{Provenance:} \cite[Prop.~3.1, Thms.~4.1, 4.2, Prop.~5.3(2)]{ButlerT12}.
\end{assembled}
\begin{proof}
The mixings of a system of $n$ conformal spinors preserving momentum and helicity for every configuration are exactly $U(n)$, acting on the $n$-tuple $(Z^{1}, \ldots, Z^{n}) \in \C^{n} \otimes \CS$ as $U \otimes 1$; the automorphism group of the hierarchy $\C^{2} \otimes \C^{3}$ is $(U(2) \times U(3))/K \cong S(U(2) \times U(3))$, and that of $\C^{3} \otimes \C^{3}$ has Lie algebra $\su(3) \oplus \su(3) \oplus \uu(1)$, dimension seventeen; the tensor ledger $\Lambda^{\bullet}(\C^{2} \otimes \C^{3})$ contains no $(2,1)$. \emph{Provenance:} \cite[Prop.~3.1, Thms.~4.1, 4.2, Prop.~5.3(2)]{ButlerT12} (\ref{c:standard-model.quarks-leptons-two-level-fermion-content}, \ref{c:standard-model.global-form-quotient-determined}, \ref{c:standard-model.gauge-algebra-minimal-massive-composite}, \ref{c:standard-model.internal-symmetry-ladder}).

\end{proof}

\begin{fact}[Imported: the Fock representation and its Bogoliubov group (external ingredients)]\label{c:rungs.imported-fock-representation-bogoliubov-group-external}
The results from outside the series on which \ref{c:rungs.imported-fock-representation-bogoliubov-group} draws, as cited there: W. Fulton and J. Harris, Representation Theory: A First Course, Graduate Texts in Mathematics 129, Springer (1991)., Lecture~20; H. B. Lawson and M.-L. Michelsohn, Spin Geometry, Princeton University Press (1989)., Ch.~I, \S\S 2, 5.
\emph{Source:} W. Fulton and J. Harris, Representation Theory: A First Course, Graduate Texts in Mathematics 129, Springer (1991)., Lecture~20; H. B. Lawson and M.-L. Michelsohn, Spin Geometry, Princeton University Press (1989)., Ch.~I, \S\S 2, 5.
\end{fact}

\begin{assembled}[Imported: the Fock representation and its Bogoliubov group]\label{c:rungs.imported-fock-representation-bogoliubov-group}
For a complex Hermitian space $E$ of dimension $n$, the Fock space $\Lambda^{\bullet}E$ is the irreducible spinor module of the complex Clifford algebra on $2n$ generators, its parity grading the chirality decomposition; the orthogonal transformations of the $2n$ Majorana generators preserving the canonical anticommutation relations form the compact $O(2n)$, implemented by $\mathrm{Pin}(2n)$, with $SO(2n)$ implemented by $\mathrm{Spin}(2n)$; those preserving the complex structure form $U(n)$, acting as $\Lambda^{\bullet}(g)$ up to a determinant phase. \emph{Provenance:} \cite[Fact~2.11]{ButlerT12}; \cite[Lecture~20]{FultonHarris1991}; \cite[Ch.~I, \S\S 2, 5]{LawsonMichelsohn1989}.
\end{assembled}
\begin{proof}
For a complex Hermitian space $E$ of dimension $n$, the Fock space $\Lambda^{\bullet}E$ is the irreducible spinor module of the complex Clifford algebra on $2n$ generators, its parity grading the chirality decomposition; the orthogonal transformations of the $2n$ Majorana generators preserving the canonical anticommutation relations form the compact $O(2n)$, implemented by $\mathrm{Pin}(2n)$, with $SO(2n)$ implemented by $\mathrm{Spin}(2n)$; those preserving the complex structure form $U(n)$, acting as $\Lambda^{\bullet}(g)$ up to a determinant phase. \emph{Provenance:} \cite[Fact~2.11]{ButlerT12} (\ref{c:standard-model.imported-fock-representation-clifford-algebra-bogoliubov}); \cite[Lecture~20]{FultonHarris1991}; \cite[Ch.~I, \S\S 2, 5]{LawsonMichelsohn1989}.

This restates \ref{c:rungs.imported-fock-representation-bogoliubov-group-external}.

\end{proof}

\begin{assembled}[Imported: results of T12 used in the registration]\label{c:rungs.imported-results-t12-used-registration}
The slot factor of the mass bracket transforms as $(1,1)_{1}$, unconditionally, and under (S) the state at $(2,0)$ is that same one-dimensional representation $\Lambda^{2}\C^{2}$. Given the content, with the fourth singlet at hypercharge zero or absent, the hypercharge assignment satisfying the anomaly conditions is unique; with a fourth singlet of free hypercharge a one-parameter family $Y + t(B - L)$ appears. Over the derived generation, which contains the lepton doublet, the trace ratio is $\sin^{2}\theta_{W} = 3/8$ (over the quark doublet and the singlets alone it is $9/26$), and the one-loop run with the measured couplings and the observed spectrum does not meet: the desert is dead and the running between the scales must change, through partner-like content or an intermediate symmetry. Under (S), in the untwisted action every abelian generator in the span of the two mode numbers except hypercharge fails three of the four anomaly conditions over the derived generation; in the action the state space carries, twisted by $\det^{-1/2}$, every abelian generator is anomaly-free, and the thirteenth generator of the five-mode ledger's automorphism algebra is the $U(1)_{X}$ of $\mathrm{Spin}(10)$, broken with the Pati-Salam algebra at the intermediate scale by the Majorana term of $\nu^{c}$, which leaves the determinant-one part of $U(2) \times U(3)$ as the gauge group below that scale. \emph{Provenance:} \cite[Prop.~3.2, Prop.~5.8(1), (3), Thm.~5.5(3), Props.~6.1, 6.2, 10.7]{ButlerT12}.
\end{assembled}
\begin{proof}
The slot factor of the mass bracket transforms as $(1,1)_{1}$, unconditionally, and under (S, \ref{c:standard-model.hyp-s-slot-statistics}) the state at $(2,0)$ is that same one-dimensional representation $\Lambda^{2}\C^{2}$. Given the content, with the fourth singlet at hypercharge zero or absent, the hypercharge assignment satisfying the anomaly conditions is unique; with a fourth singlet of free hypercharge a one-parameter family $Y + t(B - L)$ appears. Over the derived generation, which contains the lepton doublet, the trace ratio is $\sin^{2}\theta_{W} = 3/8$ (over the quark doublet and the singlets alone it is $9/26$), and the one-loop run with the measured couplings and the observed spectrum does not meet: the desert is dead and the running between the scales must change, through partner-like content or an intermediate symmetry. Under (S), in the untwisted action every abelian generator in the span of the two mode numbers except hypercharge fails three of the four anomaly conditions over the derived generation; in the action the state space carries, twisted by $\det^{-1/2}$, every abelian generator is anomaly-free, and the thirteenth generator of the five-mode ledger's automorphism algebra is the $U(1)_{X}$ of $\mathrm{Spin}(10)$, broken with the Pati-Salam algebra at the intermediate scale by the Majorana term of $\nu^{c}$, which leaves the determinant-one part of $U(2) \times U(3)$ as the gauge group below that scale. \emph{Provenance:} \cite[Prop.~3.2, Prop.~5.8(1), (3), Thm.~5.5(3), Props.~6.1, 6.2, 10.7]{ButlerT12} (\ref{c:standard-model.mass-bracket-charged-singlet}, \ref{c:standard-model.spectrum-run}, \ref{c:standard-model.hypercharge-global-form-pair-counting-consistency}, \ref{c:standard-model.under-s-composites-state-space-spinor}, \ref{c:standard-model.majorana-element}, \ref{c:standard-model.weak-angle-over-derived-generation}).

\end{proof}

\section*{What we assume}
The results below are taken as given in this paper; each result that rests on one says so.
\begin{postulate}[Hypothesis (R): the rung reading]\label{c:rungs.hyp-r-rung-reading}
\emph{Hypothesis (R), the rung reading.} The internal state space of a composite of $n$ conformal spinors (an $n$-twistor composite) is the spinor module of the $n$-th rung of the tower above $\Cl(2,4)$, $\Lambda^{\bullet}(W \oplus R_{n})$, on which the transported algebra acts as in Fact \ref{c:rungs.imported-tower-one-qubit-per-rung}: each conformal spinor of the composite occupies one new rung, and the constituents are the states with one rung occupied.
\end{postulate}

\section{The exact part: what the tower's spinor module is, and how the paper's objects sit in it}

\begin{theorem}[The $n$-rung spinor module and its mixings]\label{c:rungs.rung-spinor-module-mixings}
Let $\Cl(2{+}n, 4{+}n)$ be the $n$-th rung of the tower above $\Cl(2,4)$, with new isotropic directions $v_{1}, \ldots, v_{n}$ spanning $R_{n} \cong \C^{n}$ after complexification. Then:
\begin{enumerate}
\item Its complexified spinor module is $\Lambda^{\bullet}(W \oplus R_{n}) = \Sp \otimes \Lambda^{\bullet}\C^{n}$, of dimension $2^{3+n}$, graded by $(j, k)$ with $j$ the number of new modes and $k$ the number of old pairs occupied, and irreducible; its chirality is the parity of $j + k$. At $n = 2$ this is the five-mode ledger of Hypothesis (S), with the same grading and the same chirality rule.
\item The subgroup of unitaries of $W \oplus R_{n}$ preserving the two summands is $U(3) \times U(n)$, and it acts on the module as $\Lambda^{\bullet}(B \oplus A)$: on $\Lambda^{j}R_{n} \otimes \Lambda^{k}W$ by $\Lambda^{j}(A) \otimes \Lambda^{k}(B)$. At $n = 2$ this is exactly the action Hypothesis (S) prescribes.
\item The Bogoliubov group of the $3 + n$ modes is the compact orthogonal group $O(6{+}2n)$, implemented by $\mathrm{Pin}(6{+}2n)$. At $n = 2$ the connected part is $\mathrm{Spin}(10)$, acting irreducibly on each chirality half; the half of even $j + k$ is the sixteen-dimensional half-spin representation, which under (S) is the generation of Fact \ref{c:standard-model.under-s-composites-state-space-spinor}. The number-preserving subgroup of $\mathrm{Spin}(6{+}2n)$ is a double cover of $U(3{+}n)$, acting as $\Lambda^{\bullet} \otimes \det^{-1/2}$; on the determinant-one subgroup this agrees with the action of (2), and it is the determinant-one part that (S) uses (Fact \ref{c:rungs.imported-results-t12-used-registration}). At $n = 3$ the group is $\mathrm{Spin}(12)$, and the split-and-number-preserving $U(3) \times U(3)$, of dimension eighteen, acts faithfully on the module, while its action on the cell sector $\Lambda^{1}R_{3} \otimes \Lambda^{1}W = \C^{3} \otimes \C^{3}$ has kernel $K = \{(\lambda I_{3}, \lambda^{-1}I_{3})\}$ and image the seventeen of Fact \ref{c:rungs.imported-ladder-automorphism-group-no-lepton}: the seventeen is the $n = 3$ cell action, exactly as the twelve of $S(U(2) \times U(3))$ is the $n = 2$ cell action of a thirteen-dimensional group on the module. (At $n = 3$ the trace-free subalgebra $\{\tr X + \tr X' = 0\}$ contains $K$'s generator, so the seventeen is the quotient and not that subalgebra, unlike $n = 2$ where $\mathfrak{s}(\uu(2) \oplus \uu(3))$ maps isomorphically onto the cell action. This clause is a relabelling: the cell sector is $\C^{3} \otimes \C^{3}$ with the group Fact \ref{c:rungs.imported-ladder-automorphism-group-no-lepton} already has, and nothing about the tower enters.)
\item The dimension counts: $2^{3+n} = 8, 16, 32, 64$ for $n = 0, 1, 2, 3$; $\dim \mathfrak{so}(6{+}2n) = 15, 28, 45, 66$; $\dim(\uu(3) \oplus \uu(n)) = 9, 10, 13, 18$.
\end{enumerate}
\end{theorem}
\begin{proof}
(1) By Fact \ref{c:rungs.imported-tower-one-qubit-per-rung}, $T_{n} \cong \Cl(2,4) \otimes \mathrm{M}_{2^{n}}(\R)$, and after complexification $\Cl(2,4) \otimes \C \cong \mathrm{M}_{8}(\C)$ acting irreducibly on $\Sp = \Lambda^{\bullet}W$ (Fact \ref{c:rungs.imported-spinor-identification-witt-data-line}) while $B_{n} \otimes \C \cong \mathrm{M}_{2^{n}}(\C)$ acts irreducibly on $\Lambda^{\bullet}W_{n} \otimes \C = \Lambda^{\bullet}R_{n}$ (Fact \ref{c:rungs.imported-tower-one-qubit-per-rung}); the tensor product $\mathrm{M}_{8} \otimes \mathrm{M}_{2^{n}} = \mathrm{M}_{8 \cdot 2^{n}}$ acts irreducibly on $\Sp \otimes \Lambda^{\bullet}R_{n}$, which is the whole endomorphism algebra of that space, so this is the spinor module of $T_{n} \otimes \C$, of dimension $2^{3+n}$. The step from the twisted presentation to the mode grading: the tensor decomposition $A_{n} \otimes B_{n}$ of Fact \ref{c:rungs.imported-tower-one-qubit-per-rung} presents $T_{n}$ through the twisted ground generators, which are not vectors of $\R^{2+n,4+n}$; but $T_{n} \otimes \C$ is also generated by the $6 + 2n$ untwisted vectors, so it is the complex Clifford algebra on $6 + 2n$ generators (\ref{c:std-algebra.complexification}), $\mathrm{M}_{2^{3+n}}(\C)$, whose irreducible module is unique up to isomorphism and is the Fock module of any $3 + n$ anticommuting mode pairs (Fact \ref{c:rungs.imported-fock-representation-bogoliubov-group}); the pairs of $W$ with its dual and the untwisted $(v_{j}, w_{j})$ are such a set, and their Fock module is $\Lambda^{\bullet}(W \oplus R_{n})$ by the direct-sum formula for exterior algebras \cite[App.~B.2]{FultonHarris1991}, with the grading by occupied modes the bigrading $\Lambda^{j}R_{n} \otimes \Lambda^{k}W$. The irreducibility established through $A_{n} \otimes B_{n}$ transfers to this presentation because the module is the same up to isomorphism. The chirality of a Fock module is its parity grading (Fact \ref{c:rungs.imported-fock-representation-bogoliubov-group}), here $(-1)^{j+k}$. At $n = 2$, $R_{2} = \C^{2}$ and the description is (S)'s. (2) A unitary of $W \oplus R_{n}$ preserving the summands is $B \oplus A$, and the induced action on the exterior algebra is $\Lambda^{\bullet}(B \oplus A) = \Lambda^{\bullet}(B) \otimes \Lambda^{\bullet}(A)$ on $\Lambda^{\bullet}W \otimes \Lambda^{\bullet}R_{n}$, degree by degree. (3) is Fact \ref{c:rungs.imported-fock-representation-bogoliubov-group} with $E = W \oplus R_{n}$, $\dim E = 3 + n$, the Hermitian structure being the one under which $v_{j}^{*} = w_{j}$ (Fact \ref{c:rungs.imported-tower-one-qubit-per-rung}); the double cover and the $\det^{-1/2}$ twist are the "determinant phase" of Fact \ref{c:rungs.imported-fock-representation-bogoliubov-group}, absent on the determinant-one subgroup. The $n = 3$ clause: $\uu(3) \oplus \uu(3)$ has dimension $18$; $K$'s generator $(iI_{3}, -iI_{3})$ acts on $\Lambda^{j}R_{3} \otimes \Lambda^{k}W$ by $i(j - k)$, zero exactly on the cells $j = k = 1$ and nonzero on the module, so the action on the module is faithful and the action on the cells has kernel $K$; the quotient, of dimension $17$, is the seventeen of Fact \ref{c:rungs.imported-ladder-automorphism-group-no-lepton}; and the trace-free subalgebra, which contains the kernel since $\tr(iI_{3}) + \tr(-iI_{3}) = 0$, acts on the cells with rank $16$ only. The script confirms the ranks $17$ and $16$ at $n = 3$ against $12$ and $12$ at $n = 2$, and the eigenvalue $j - k$. (4) is arithmetic, script-checked.

\end{proof}

\begin{remark}[The compact form]\label{c:note.rungs.compact-form}
The ambient signature at rung two is $(4,6)$, whose spin group is the noncompact $\mathrm{Spin}(4,6)$; the group in Theorem \ref{c:rungs.rung-spinor-module-mixings}(3) is the compact $\mathrm{Spin}(10)$. They are two real forms of one complex group sharing the module, and the compact one is the real form fixed by the $*$-structure of the tower (Fact \ref{c:rungs.imported-tower-one-qubit-per-rung}), under which the ladder pairs are mutually adjoint, $v_{j}^{*} = w_{j}$: that adjointness is the Hermitian structure of Fact \ref{c:rungs.imported-fock-representation-bogoliubov-group}, and the Bogoliubov group preserving it is the compact one. The $*$-structure is defined once the tower's frames are fixed, as T3 fixes them, so "fixed by" means fixed given the frames, a choice of maximal compact subgroup made once for the whole series rather than a consequence of the relations alone. The pattern is already present one rung down, where the ambient real form is the conformal $SU(2,2) \cong \mathrm{Spin}(2,4)$ (Fact \ref{c:rungs.imported-spinor-identification-witt-data-line}) while the pair-mixing group is the unitary group $U(3)$ of $W$ for the same adjointness, inside the compact three-mode Bogoliubov group. A grand-unification group must be compact, and the tower supplies the compact form for the same reason it supplies $U(3)$.
\end{remark}

\begin{lemma}[The mass bracket is the frozen point's grip on the slot-antisymmetric bivector, at $\Lambda = 0$]\label{c:rungs.mass-bracket-frozen-points-grip-slot}
In the chart $\CS = \Som \oplus \Spi$ of the conformal spinor space (twistor space), the $E = -1$ and $E = +1$ eigenspaces, with bases $e_{1}, e_{2}$ of $\Som$ and $\pi_{1}, \pi_{2}$ of $\Spi$, normalise $\varepsilon(e_{1} \wedge e_{2} \wedge \pi_{1} \wedge \pi_{2}) = 1$ and let $n_{\infty} = e_{1} \wedge e_{2} \in \Lambda^{2}\Som \subset \Lambda^{2}\CS \cong V_{\C}$, the bivector of the line at infinity. For conformal spinors $Z^{i} = (\omega^{i}, \pi^{i})$,
\[
\varepsilon\bigl(n_{\infty} \wedge Z^{1} \wedge Z^{2}\bigr) \;=\; \br{\pi^{1}\pi^{2}},
\qquad\text{so}\qquad
m^{2} \;=\; 2\,\bigl|\varepsilon(n_{\infty} \wedge Z^{1} \wedge Z^{2})\bigr|^{2}
\]
by Fact \ref{c:rungs.imported-mass-symplectic-area}; and $\varepsilon(n_{0} \wedge Z^{1} \wedge Z^{2}) = \br{\omega^{1}\omega^{2}}$ for $n_{0} = \pi_{1} \wedge \pi_{2}$, so that the frozen point at general $\Lambda$ pairs to $\varepsilon(I_{\Lambda} \wedge Z^{1} \wedge Z^{2}) = \br{\pi^{1}\pi^{2}} + (\Lambda/6)\br{\omega^{1}\omega^{2}}$, the $\Lambda = 0$ term being the mass bracket. In the language of the $n$-tuple $\Psi = \sum_{i} e_{i}' \otimes Z^{i} \in \C^{2} \otimes \CS$ (with $e_{i}'$ the slot basis, over the complex unit $I_{6}$), the total momentum is the slot trace of $\Psi \otimes \bar\Psi$ restricted to the $\pi$-parts, and $Z^{1} \wedge Z^{2}$ is the $\Lambda^{2}\C^{2} \otimes \Lambda^{2}\CS$ component of $\Psi \odot \Psi$: at $\Lambda = 0$ the mass is the frozen point's pairing with the slot-antisymmetric bivector, and vanishes exactly when that bivector is orthogonal to the frozen point.
\end{lemma}
\begin{proof}
Expand $Z^{1} \wedge Z^{2} = \omega^{1}\wedge\omega^{2} + (\omega^{1}\wedge\pi^{2} - \omega^{2}\wedge\pi^{1}) + \pi^{1}\wedge\pi^{2}$. Wedging with $e_{1} \wedge e_{2}$ kills every term with an $\omega$ factor, leaving $e_{1}\wedge e_{2}\wedge\pi^{1}\wedge\pi^{2} = \det(\pi^{i}_{A})\,e_{1}\wedge e_{2}\wedge\pi_{1}\wedge\pi_{2}$, and $\det(\pi^{i}_{A}) = \br{\pi^{1}\pi^{2}}$; the $n_{0}$ case is the same with the roles exchanged. That $n_{\infty}$ is the bivector of $\Som$ and is the frozen point at $\Lambda = 0$ is Facts \ref{c:rungs.imported-spinor-identification-witt-data-line} and \ref{c:rungs.imported-mass-symplectic-area}. For the tuple: $\Psi \otimes \bar\Psi = \sum_{ij} e_{i}' \otimes \bar e_{j}' \otimes Z^{i} \otimes \bar Z^{j}$, and contracting the slot indices with $\delta_{ij}$ gives $\sum_{i} Z^{i} \otimes \bar Z^{i}$, whose $\pi \otimes \bar\pi$ part is $p$; and $S^{2}(\C^{2} \otimes \CS) = S^{2}\C^{2} \otimes S^{2}\CS \oplus \Lambda^{2}\C^{2} \otimes \Lambda^{2}\CS$, the second summand of $\Psi \odot \Psi$ being $(e_{1}' \wedge e_{2}') \otimes (Z^{1} \wedge Z^{2})$. The script checks the two volume identities on rational instances.

\end{proof}

\begin{proposition}[The ladder group is the rung-mixing group on the one-constituent sector]\label{c:rungs.ladder-group-rung-mixing-group-one}
The $U(n)$ of Fact \ref{c:rungs.imported-ladder-automorphism-group-no-lepton} acts on the $n$-tuple $\Psi \in \C^{n} \otimes \CS$ of conformal spinors (twistors), $\CS$ the conformal spinor space, as $U \otimes 1$. Under Theorem \ref{c:rungs.rung-spinor-module-mixings}(2) the rung-mixing $U(n)$ acts on the sector $\Lambda^{1}R_{n} \otimes \Lambda^{\mathrm{even}}W = \C^{n} \otimes \CS$ as $A \otimes 1$. The two actions coincide: the ladder of the Standard-Model paper is the mixing of the new rungs of the tower, restricted to the states with exactly one new mode occupied and an even number of old pairs. The same holds for the pair-mixing $U(3)$, which acts on $\CS$ by $\Lambda^{\mathrm{even}}(B)$ in both.
\end{proposition}
\begin{proof}
Both are $A \otimes 1$ on $\C^{n} \otimes \CS$ by construction, Fact \ref{c:rungs.imported-ladder-automorphism-group-no-lepton} for the first and Theorem \ref{c:rungs.rung-spinor-module-mixings}(2) at $j = 1$, $k$ even, for the second; the identification $\Lambda^{1}R_{n} = \C^{n}$ sends the slot basis to the new modes $v_{1}, \ldots, v_{n}$. The $U(3)$ statement is Theorem \ref{c:rungs.rung-spinor-module-mixings}(2) at $j = 1$ read on the $W$ factor. The proposition is definitional once the slot basis is identified with the rung basis; its content is that this identification is available, not that it is forced, and the script computes the rung operators of $\uu(2)$ on the sector and finds the ladder action and nothing leaving the sector.

\end{proof}

\begin{remark}[What this identifies, and what it does not]\label{c:note.rungs.what-identifies-what-it-does-not}
Proposition \ref{c:rungs.ladder-group-rung-mixing-group-one} is exact, and it is the one place where the tower and the Standard-Model paper's objects are literally the same thing: the ladder's $U(n)$, derived from momentum and helicity preservation, and the tower's rung-mixing group, derived from the completion, agree on $\C^{n} \otimes \CS$, $\CS$ the conformal spinor space (twistor space). But the identification is on that sector only. The ladder theorem says nothing about the other sectors ($j \neq 1$, or $k$ odd), and the tower says nothing about momentum. To say that the composite's state space is the whole module is a hypothesis, stated next.
\end{remark}

\phantomsection\label{c:note.rungs.setting-exact-part-what-towers-spinor}
\begin{quote}
\emph{Hypothesis (R), the rung reading.} The internal state space of an $n$-twistor composite is the spinor module of the $n$-th rung of the tower above $\Cl(2,4)$, $\Lambda^{\bullet}(W \oplus R_{n})$, on which the transported algebra acts as in Fact \ref{c:rungs.imported-tower-one-qubit-per-rung}: each twistor of the composite occupies one new rung, and the constituents are the states with one rung occupied.
\end{quote}

\begin{proposition}[(R) and (S) are one hypothesis in two languages]\label{c:rungs.r-s-are-one-hypothesis-two}
Given Facts \ref{c:rungs.imported-tower-one-qubit-per-rung} and \ref{c:rungs.imported-fock-representation-bogoliubov-group}, Hypothesis (R) at $n = 2$ implies Hypothesis (S); and (S), together with the identification of the two slots with the two rungs, implies (R) at $n = 2$. At $n = 2$ that identification is a relabelling of two basis vectors, so (R) and (S) are equivalent up to it. Under either, the constituent hypothesis of T12's successor problem (6), a slot vector times a Dirac state with the slot mixing acting once, holds as the $j = 1$ sector; the derived generation is the even half of the module (Fact \ref{c:standard-model.under-s-composites-state-space-spinor}); the mass bracket is the pairing of Lemma \ref{c:rungs.mass-bracket-frozen-points-grip-slot}; and $\mathrm{Spin}(10)$ is the Bogoliubov group of the composite's modes.
\end{proposition}
\begin{proof}
(R) $\Rightarrow$ (S): by Theorem \ref{c:rungs.rung-spinor-module-mixings}(1) and (2) the module of (R) is $\Lambda^{\bullet}(W \oplus \C^{2})$ with the $(j,k)$ grading, the action $\Lambda^{\bullet}(A \oplus B)$ of the rung-and-pair mixing group and chirality $(-1)^{j+k}$. (S) speaks of the automorphism group of the hierarchy (Fact \ref{c:rungs.imported-ladder-automorphism-group-no-lepton}) acting by $\Lambda^{\bullet}(A \oplus B)$ (the $U(n)$ action on the Fock space of \ref{c:rungs.imported-fock-representation-bogoliubov-group}); the two groups are identified on the one-constituent sector by Proposition \ref{c:rungs.ladder-group-rung-mixing-group-one}, and (R)'s clause that each conformal spinor occupies one rung extends the identification of slots with rungs, hence of the hierarchy's $U(2) \times U(3)$ with the rung-and-pair mixing group, to the whole module. With that identification the module, grading, action and chirality are (S) clause by clause. (S) $\Rightarrow$ (R): (S)'s space is $\Lambda^{\bullet}(\C^{2} \oplus W)$, which by Theorem \ref{c:rungs.rung-spinor-module-mixings}(1) is the two-rung module, and (S)'s action is the rung-and-pair mixing of Theorem \ref{c:rungs.rung-spinor-module-mixings}(2); the transported algebra's action is the Chevalley action on the exterior algebra (Fact \ref{c:rungs.imported-tower-one-qubit-per-rung}), which (S)'s space carries, so (R) holds once slots are identified with rungs, which is the identification the statement names. The remaining clauses are Fact \ref{c:standard-model.under-s-composites-state-space-spinor}, Lemma \ref{c:rungs.mass-bracket-frozen-points-grip-slot}, and Theorem \ref{c:rungs.rung-spinor-module-mixings}(3). The proposition is a relabelling made exact, and is stated as a proposition rather than a theorem for that reason: its content is that the two hypotheses have the same module, grading, group action and chirality, so that nothing is gained or lost in passing between them, and that the passage adds exactly one identification, slots with rungs.

\end{proof}

\begin{proposition}[The generators outside the gauge algebra]\label{c:rungs.generators-outside-gauge-algebra}
Under (R) or (S), of the forty-five generators of $\mathfrak{so}(10)$ acting on the composite's modes, exactly thirteen preserve the slot-pair split and the mode number, spanning $\uu(2) \oplus \uu(3)$; exactly twenty-five preserve the mode number, spanning $\uu(5)$; so thirty-two lie outside $\uu(2) \oplus \uu(3)$, twelve of them inside $\uu(5)$, moving one excitation between a rung and a pair, and twenty changing the mode number by two.
\end{proposition}
\begin{proof}
Here the imaginary unit is $I_{4}$, which commutes with every mode operator. The Majorana generators $\gamma_{2i} = a_{i} + a_{i}^{\dagger}$, $\gamma_{2i+1} = I_{4}(a_{i} - a_{i}^{\dagger})$ of the five modes (by \ref{c:rungs.hyp-r-rung-reading} or \ref{c:standard-model.hyp-s-slot-statistics}) satisfy the Clifford relations, and the forty-five bilinears $\gamma_{p}\gamma_{q}$, $p < q$, span $\mathfrak{so}(10)$ (Fact \ref{c:rungs.imported-fock-representation-bogoliubov-group}). The bilinears from one mode, $\gamma_{2i}\gamma_{2i+1} = I_{4}(2a_{i}^{\dagger}a_{i} - 1)$, five of them, and from two modes $i \neq j$ the two combinations $a_{i}^{\dagger}a_{j} - a_{j}^{\dagger}a_{i}$ and $I_{4}(a_{i}^{\dagger}a_{j} + a_{j}^{\dagger}a_{i})$, twenty of them, commute with the total number operator and span $\uu(5)$, of dimension twenty-five; the remaining two combinations from each pair of modes, $a_{i}a_{j} \pm a_{i}^{\dagger}a_{j}^{\dagger}$ up to factors of $I_{4}$, twenty of them, change the mode number by two. Of the twenty-five, those with $i, j$ both rungs or both pairs commute with the rung and pair numbers separately: $5 + 2 + 6 = 13$, spanning $\uu(2) \oplus \uu(3)$; the other twelve, with one rung index and one pair index, move an excitation across the split. The script computes the two centralisers as exact kernels, of dimensions $25$ and $13$, and exhibits the spanning sets.

\end{proof}

\section{The part that does not close: no mechanism makes a twistor a rung}

\phantomsection\label{c:note.rungs.setting-part-does-not-close-no}
Proposition \ref{c:rungs.r-s-are-one-hypothesis-two} replaces one hypothesis by another of the same strength, and the gain is only that (R) is a sentence about the tower, an object the series has theorems about, rather than a sentence about slot statistics, which it does not. A derivation of (R) would need a mechanism producing one rung per twistor. Three candidates present themselves, and none works as stated.

\begin{corollary}[Null momenta invite one rung, together]\label{c:rungs.fact-null-momenta-invite-one-rung}
Let $p^{1}, \ldots, p^{n}$ be null vectors of $\R^{1,3} \subset \Cl(1,3)$, the momenta of $n$ massless conformal spinors (twistors). No extension of $\Cl(1,3)$ carries a $\mathrm{Spin}^{0}$-equivariant element-level inverse of the null vectors, so none inverts the $p^{i}$ in the sense the transport principle asks; the minimal completion that the transport principle prescribes is the single rung $\Cl(2,4)$, unique up to isometry and independent of $n$; and on that completion's quadric the transported reciprocal is total, sending every $p^{i}$ to its ideal point $[p^{i} + \tfrac12 n_{\infty}]$ at once. The transport principle completes the quadratic space, not each zero divisor: a second null momentum invites no second rung.
\end{corollary}
\begin{proof}
All three clauses are Fact \ref{c:rungs.imported-totality-minimality-completion} applied to $(\R^{1,3}, \eta)$, which is nondegenerate of mixed signature: the no-go for element-level inverses, the minimality and uniqueness of the one-plane completion, and the totality of the reciprocal on the quadric. The rung is determined by the quadratic space alone, whose null cone contains every $p^{i}$, and the reciprocal's totality is a statement about all null vectors of that cone simultaneously. A second conformal spinor changes the configuration, not the quadratic space. The corollary adds nothing to Fact \ref{c:rungs.imported-totality-minimality-completion} beyond the observation that its hypotheses concern the quadratic space and not the number of null vectors named. The script confirms, in the blade engine, that a null $p$ stays null in $\Cl(2,4)$ and in $\Cl(3,5)$ (the rung is an isometric extension and inverts nothing as an element) and that the rungs supply one further anticommuting null pair each, orthogonal to $p$.

\end{proof}

\begin{proposition}[The bivector of two conformal spinors is null, invites one rung, and that rung's $U(2)$ mixes chirality]\label{c:rungs.two-twistor-bivector-null-invites-one}
For conformal spinors (twistors) $Z^{1}, \ldots, Z^{n} \in \CS$, where $\CS = \Lambda^{\mathrm{even}}W$ is the conformal spinor space with complex unit $I_6$ (not the Dirac module $\Sp = \Lambda^{\bullet}W$ below), every pairwise bivector $Z^{i} \wedge Z^{j} \in \Lambda^{2}\CS \cong V_{\C}$ is a null vector of the complexified ambient space of $\Cl(2,4)$, decomposable by construction, and for $n = 2$, $m^{2} = 2|\varepsilon(n_{\infty} \wedge Z^{1}\wedge Z^{2})|^{2}$ by Lemma \ref{c:rungs.mass-bracket-frozen-points-grip-slot}. The transport principle iterates on the real ambient space $\R^{2,4}$ (Fact \ref{c:rungs.imported-tower-one-qubit-per-rung}), and Fact \ref{c:rungs.imported-totality-minimality-completion} applied to it gives one further hyperbolic plane, unique, with a reciprocal total on the whole real cone, whatever $n$: the mechanism is a completion of one quadratic space and depends on no count of conformal spinors. The one-rung module $\Lambda^{\bullet}(W \oplus \C v)$, of dimension sixteen, is $\C^{2} \otimes \Sp$ as a vector space, the constituent space of (C), and the transported algebra does supply a $U(2)$ on the $\C^{2}$ factor: the unitary group of the rung's own Clifford algebra $B_{1} \otimes \C \cong \mathrm{M}_{2}(\C)$, the commutant of the ground algebra. But that $U(2)$ contains the odd elements $v \pm w$ of the rung, which anticommute with the chirality operator; its chirality-preserving part is the abelian $\{1, vw\}$-span, a torus. So one rung gives (C)'s vector space with a $U(2)$ that mixes the chirality of the two constituent states, while (C), (S) and the ladder need a $U(2)$ acting within one chirality, which is the rung-mixing $U(2)$ on $\Lambda^{1}R_{2}$ of Theorem \ref{c:rungs.rung-spinor-module-mixings}(2), available at two rungs.
\end{proposition}
\begin{proof}
Decomposability is the form of $Z^{i} \wedge Z^{j}$; nullity of decomposable bivectors is Fact \ref{c:rungs.imported-spinor-identification-witt-data-line}; the mass identity is Lemma \ref{c:rungs.mass-bracket-frozen-points-grip-slot}. Whether a given complex null bivector lies on the real cone is the nullity of its line for the reversion form $\Sigma$, the twistor form (Fact \ref{c:rungs.imported-spinor-identification-witt-data-line}), which a massive pair need not satisfy; the argument does not need it, since it is the real quadratic space $\R^{2,4}$ that is completed, and $\R^{2,4}$ is nondegenerate of mixed signature, so Fact \ref{c:rungs.imported-totality-minimality-completion} applies to it exactly as to $\R^{1,3}$ in Corollary \ref{c:rungs.fact-null-momenta-invite-one-rung}. One rung above $\Cl(2,4)$ has module $\Lambda^{\bullet}W \otimes \Lambda^{\bullet}(\C v) = \Sp \otimes \C^{2}$ and $T_{1} \cong A_{1} \otimes B_{1}$ with $B_{1} \cong \Cl(1,1) \cong \mathrm{M}_{2}(\R)$ acting irreducibly on the qubit and commuting with the twisted ground algebra $A_{1}$ (Fact \ref{c:rungs.imported-tower-one-qubit-per-rung}); after complexification $A_{1} \otimes \C \cong \mathrm{M}_{8}(\C)$, so the commutant of $A_{1}$ in $\mathrm{End}(\Sp \otimes \C^{2}) = \mathrm{M}_{8} \otimes \mathrm{M}_{2}$ is $1 \otimes \mathrm{M}_{2}(\C)$, whose unitary group is $U(2)$. The rung's vectors $v, w$ are odd for the mode-number grading, so they anticommute with $(-1)^{N}$, and the even part of $\mathrm{M}_{2}(\C)$ is spanned by $1$ and $vw$, which commute. At two rungs the new modes span $\C^{2}$ and $U(2) \subset U(5)$ acts on $\Lambda^{1}R_{2} = \C^{2}$ within a fixed chirality (Theorem \ref{c:rungs.rung-spinor-module-mixings}(2)). The script checks the nullity of all six pairwise bivectors of an instance with four conformal spinors, and, on the four-mode Fock space, that the twisted ground generators generate $\mathrm{M}_{8}(\C)$, that the rung's $\{1, v, w, vw\}$ commutes with them, that $v$ and $w$ anticommute with chirality, and that the even part is abelian.

\end{proof}

\begin{proposition}[The transported module holds the tuple as one vector, not as two conformal spinors]\label{c:rungs.transported-module-holds-tuple-one-vector}
The configuration space of two conformal spinors (two twistors) of Fact \ref{c:rungs.imported-mass-symplectic-area}, as a vector space, is $\CS \oplus \CS = \C^{2} \otimes \CS$, $\CS$ the conformal spinor space, eight-dimensional, carrying the continuous data $(\omega^{i}, \pi^{i})$, the parts of each spinor in the $E = -1$ and $E = +1$ eigenspaces at the origin; it is the sector $(j, k) \in \{(1,0), (1,2)\}$ of the two-rung module, which under $U(2) \times U(3)$ is, by Fact \ref{c:standard-model.under-s-composites-state-space-spinor}, $(2,1)_{1/2} \oplus (2,\bar{3})_{-1/6}$, the antiparticle doublets $L^{*} \oplus Q^{*}$ in the orientation of T12's listing or $L \oplus Q$ in the opposite one. The two-particle tensor product $\CS \otimes \CS$, by contrast, has no $U(3)$-equivariant embedding in the module: it contains $S^{2}(\Lambda^{2}W)$, the six-dimensional $SU(3)$-irreducible $\bar{6}$, while every sector of the module is a multiple of $\Lambda^{\bullet}W$, whose $SU(3)$ content is $1, 3, \bar{3}, 1$. And the mass bracket, the $\Lambda^{2}\C^{2} \otimes \Lambda^{2}\CS$ component of $\Psi \odot \Psi$ (Lemma \ref{c:rungs.mass-bracket-frozen-points-grip-slot}), is quadratic in $\Psi$ and lives in $S^{2}$ of the module, not in the module. So the module treats the composite as T12 does, as one tuple $\Psi$ with a slot index, and not as a product state of two conformal spinors.
\end{proposition}
\begin{proof}
The sector statement is Theorem \ref{c:rungs.rung-spinor-module-mixings}(2) at $j = 1$, $k \in \{0, 2\}$, with the weights $3 - 2k = 3, -1$ of Fact \ref{c:standard-model.under-s-composites-state-space-spinor}, and $\Lambda^{1}R_{2} \otimes \Lambda^{\mathrm{even}}W = \C^{2} \otimes \CS$ as in Proposition \ref{c:rungs.ladder-group-rung-mixing-group-one}. The $\bar{6}$: $\Lambda^{2}W \otimes \Lambda^{2}W = S^{2}(\Lambda^{2}W) \oplus \Lambda^{2}(\Lambda^{2}W)$, and $S^{2}(\Lambda^{2}W)$ has the weight $2(e_{1} + e_{2})$, whose class modulo $(1,1,1)$ is not a weight of $\Lambda^{\bullet}W$, all of whose weights have coordinates in $\{0, 1\}$; so $S^{2}(\Lambda^{2}W)$, irreducible under $SU(3)$ \cite[\S 15.3]{FultonHarris1991}, has no equivariant image in $\Lambda^{\bullet}W \otimes \Lambda^{\bullet}R_{2}$, on which $SU(3)$ acts through $\Lambda^{\bullet}W$ alone, and an equivariant map from $\CS \otimes \CS$ must kill it, so is not an embedding. The bracket clause is Lemma \ref{c:rungs.mass-bracket-frozen-points-grip-slot}. The script checks the weights, the Casimirs of the sectors, and the absent $\bar{6}$ weights.

\end{proof}

\begin{remark}[Where this leaves (R)]\label{c:note.rungs.where-leaves-r}
Three things are now exact that were not. The five-mode ledger is the tower's two-rung module, so (S) is a statement about the tower. The ladder group is the rung-mixing group on the one-constituent sector. And the mass bracket is the frozen point's grip on the slot-antisymmetric bivector, at $\Lambda = 0$, which is a null vector of the ambient space that invites the next rung. What is missing is a single mechanism: why a conformal spinor constituent (a twistor) occupies a rung. The count is the obstacle. Null momenta invite one rung in total (Corollary \ref{c:rungs.fact-null-momenta-invite-one-rung}); the bivectors of two conformal spinors invite one more, together (Proposition \ref{c:rungs.two-twistor-bivector-null-invites-one}); neither gives one rung per constituent, and the second gives the right vector space with a $U(2)$ that mixes chirality. The precise form of the missing statement is this: the composite's internal algebra is generated by $3 + n$ anticommuting mode pairs, the three pairs of $W$ shared and one new pair per constituent. Fact \ref{c:rungs.imported-tower-one-qubit-per-rung} then supplies everything else, the module, the grading and the groups. Why a conformal spinor constituent should contribute exactly one mode to the composite, rather than none (the tensor ledger of Fact \ref{c:rungs.imported-ladder-automorphism-group-no-lepton}, where the slots are labels and the pairs are shared) or three (a copy of $W$ per constituent, the product state that Proposition \ref{c:rungs.transported-module-holds-tuple-one-vector} excludes), is the question, and the note has no argument for it. It belongs to the tower paper, where the mode count is a theorem about the real Witt index and where a statement about what a conformal spinor constituent contributes to a composite's algebra would have to be made.
\end{remark}

\section{The fallback: (S) as a Tier C premise}

\phantomsection\label{c:note.rungs.setting-fallback-s-tier-c-premise}
If the missing statement of Remark \ref{c:note.rungs.where-leaves-r} is not supplied, the series' own rule for a structural premise that resists derivation is to register it at Tier C, with its predictions and its kill conditions priced. The registration, in the form T12 uses for its stances:

\begin{quote}
\emph{Registered premise (S), equivalently (R).} The internal state space of the minimal massive composite is the five-mode Fock space $\Lambda^{\bullet}(W \oplus \C^{2})$, the two-rung transport of the twistor completion, with the gauge group below the intermediate scale acting through Theorem \ref{c:rungs.rung-spinor-module-mixings}(2) on its determinant-one part (the thirteenth generator, anomaly-free in the action the state space carries, being broken at that scale, Fact \ref{c:rungs.imported-results-t12-used-registration}) and chirality the parity of the mode number.
\end{quote}

\noindent\emph{What it predicts.} (i) Exactly one generation per composite, the sixteen of $\mathrm{Spin}(10)$: quark and lepton doublets, four singlets, the fourth at hypercharge zero (Fact \ref{c:standard-model.under-s-composites-state-space-spinor}); that the composite carries no further chiral state is the registration's own reading of "the state space is this Fock space", not a clause of Fact \ref{c:standard-model.under-s-composites-state-space-spinor}. (ii) The $\mathbb{Z}_{6}$ charge lattice of T12 Theorem 4.2; this is unconditional in T12 and is not a test of (S), and is listed because (S) must respect it and does (Fact \ref{c:standard-model.under-s-composites-state-space-spinor}). (iii) Unification into $\mathrm{Spin}(10)$ with $\sin^{2}\theta_{W} = 3/8$ at the unification scale, hence, with the desert dead (Fact \ref{c:rungs.imported-results-t12-used-registration}), a change in the running between the weak and unification scales: partner-like content, or an intermediate symmetry such as T12's Pati-Salam route, on which $\sin^{2}\theta_{W}$ is an input. (iv) Baryon and lepton number violation by the thirty-two generators of $\mathfrak{so}(10)$ outside $\uu(2) \oplus \uu(3)$ (Proposition \ref{c:rungs.generators-outside-gauge-algebra}): twelve that move an excitation between a rung and a pair (the leptoquark directions of $\uu(5)$), and twenty that create or annihilate two modes at once.

\noindent\emph{What kills it.} (a) A confirmed chiral fermion outside the sixteen per generation, or a fourth colour singlet with nonzero hypercharge (the anomaly conditions would then admit a family, Fact \ref{c:rungs.imported-results-t12-used-registration}). (b) A fractionally charged colour singlet, off the lattice; this kills the exclusion stance of T12 Theorem 4.2 whether or not (S) holds, so it is a kill shared with T12, not one that discriminates (S). (c) A measured $\sin^{2}\theta_{W}$ at a scale where the three couplings can be shown, with the observed spectrum plus any partner content compatible with the sixteen, not to meet: the two-loop run with thresholds, T12 successor problem (2). (d) A proton lifetime bound incompatible with any $\mathrm{Spin}(10)$ breaking scale allowed by (iii); the framework does not yet fix that scale, so this condition is priced as inherited from the $\mathrm{Spin}(10)$ literature rather than sharpened. (e) Internally: a proof that the composite's internal algebra closes on the three pair modes alone, with the slots as labels (the tensor ledger reading of Fact \ref{c:rungs.imported-ladder-automorphism-group-no-lepton}), which would refute (R) directly and leave (S) without its only structural reading.

\noindent\emph{Level.} A structural premise with priced kill conditions is Tier C by the series' definition; T12 is declared Tier B, and if (S) is registered rather than derived the sections that rest on it (T12 Proposition 5.3(2), Theorem 5.5(1), the consistency clause of (3), the last clause of Proposition 5.6, Proposition 5.8, Remark 5.9, and Proposition 6.1's values $3/8$ and $3/4$, which are taken over a generation containing the lepton; without the lepton the ratio is $9/26$) carry Tier C and the abstract should say so. The sections that do not rest on it, the algebra, the global form, the converse, the quark doublet, the singlet spectrum as a computation on the single-twistor Fock space, the uniqueness arithmetic given a charge list, the trace ratios as arithmetic over any tabulated content with $9/10$ over the fibre, and the run, stay at Tier B.

\section{Position in the series}

\phantomsection\label{c:note.rungs.setting-position-series}
This note is a companion to T12 of the series, alongside a companion note on the lepton, which asks whether the lepton is derivable inside the two-level hierarchy and answers no. \emph{It relies on:} the tower paper T3 \cite{ButlerT3} for the rungs, the Fock module and the Witt index (Fact \ref{c:rungs.imported-tower-one-qubit-per-rung}), and the transport principle paper T1 \cite{ButlerT1} for the totality and minimality of the completion (Fact \ref{c:rungs.imported-totality-minimality-completion}); the twistor paper T9 \cite{ButlerT9} for the spinor identification, the Witt data and the line at infinity (Fact \ref{c:rungs.imported-spinor-identification-witt-data-line}); the mass paper T11 \cite{ButlerT11} for the two-twistor mass formula (Fact \ref{c:rungs.imported-mass-symplectic-area}); and the Standard-Model paper T12 \cite{ButlerT12} for the ladder, the automorphism group, Hypothesis (S), the derived generation, the constituent hypothesis of its successor problem (6), and the results priced in the section ``The fallback: (S) as a Tier C premise'' (Facts \ref{c:rungs.imported-ladder-automorphism-group-no-lepton}, \ref{c:rungs.imported-fock-representation-bogoliubov-group}, \ref{c:standard-model.under-s-composites-state-space-spinor}, \ref{c:rungs.imported-results-t12-used-registration}). \emph{Papers that would rely on it:} T12, if (R) replaces (S) as the premise or is derived; the tower paper T3, if the missing statement of Remark \ref{c:note.rungs.where-leaves-r} is taken up there, where it belongs; and the capstone T17 \cite{ButlerT17}, which would record (S) as a registered premise with the kill conditions of the section ``The fallback: (S) as a Tier C premise''.

\section{Machine verification}

\phantomsection\label{c:note.rungs.setting-machine-verification}
Every computational identity asserted here is verified in exact arithmetic by the deposited script \cite{ButlerVerifyT12R}, \texttt{verify\_T12\_rungs.py}, which imports the series' shared library \texttt{GA.py}: seventeen checks, all passing, no floating point (ranks are computed by \texttt{GA.py}'s elimination over exact rationals). The script builds the five-mode Fock representation explicitly, as thirty-two by thirty-two matrices over the Gaussian integers, and the four-mode one of a single rung, and computes on them rather than restating formulas. They are: the canonical anticommutation relations, the uniqueness of the vacuum and the spanning of the module by creation monomials (Fact \ref{c:rungs.imported-fock-representation-bogoliubov-group}, Theorem \ref{c:rungs.rung-spinor-module-mixings}(1)); the chirality operator $(-1)^{N}$ as the parity of $j + k$, anticommuting with every generator; the bigraded dimensions counted on the basis; the ten Majorana generators' Clifford relations and the linear independence of the forty-five bilinears; their closure under the commutator with the $\mathfrak{so}(10)$ bracket on all pairs, and their commuting with chirality (Theorem \ref{c:rungs.rung-spinor-module-mixings}(3)); the exact dimensions $25$ and $13$ of the subalgebras commuting with the total and with the separate mode numbers, giving $45 - 13 = 32 = 12 + 20$ (Theorem \ref{c:rungs.rung-spinor-module-mixings}(2), (3), Proposition \ref{c:rungs.generators-outside-gauge-algebra}, the section ``The fallback: (S) as a Tier C premise''(iv)); the selection of $3N_{\mathrm{rung}} - 2N_{\mathrm{pair}}$ as the trace-free element of the centre, with eigenvalue $3j - 2k$, and the eigenvalue $j - k$ of $K$'s generator, zero on the cells only (Proposition \ref{c:rungs.r-s-are-one-hypothesis-two}, Theorem \ref{c:rungs.rung-spinor-module-mixings}(3)); the $\su(2)$ and $\su(3)$ Casimirs on every sector and the weights of $\Lambda^{2}W$ (Proposition \ref{c:rungs.transported-module-holds-tuple-one-vector}, Fact \ref{c:standard-model.under-s-composites-state-space-spinor}); the ranks $17$ and $16$ at $n = 3$ against $12$ and $12$ at $n = 2$ (Theorem \ref{c:rungs.rung-spinor-module-mixings}(3)); the two volume identities of Lemma \ref{c:rungs.mass-bracket-frozen-points-grip-slot} on rational instances and the partial trace of $\Psi \otimes \bar\Psi$ over the slot index on a Gaussian-rational instance; the rung operators of $\uu(2)$ on the one-constituent even sector as the ladder action (Proposition \ref{c:rungs.ladder-group-rung-mixing-group-one}); a null momentum staying null in $\Cl(2,4)$ and $\Cl(3,5)$ with the rungs' null pairs in \texttt{GA.py}'s blade engine (Corollary \ref{c:rungs.fact-null-momenta-invite-one-rung}); the nullity of all six pairwise bivectors of a four-twistor instance, and, on the four-mode space, the twisted ground generators generating $\mathrm{M}_{8}(\C)$ with the rung's $\{1, v, w, vw\}$ as commutant, $v, w$ anticommuting with chirality and the even part abelian (Proposition \ref{c:rungs.two-twistor-bivector-null-invites-one}; the script's twisted generators square to $-1$, so this verifies the complexified statement, not Fact \ref{c:rungs.imported-tower-one-qubit-per-rung}'s real isomorphism); the absent $\bar{6}$ weights (Proposition \ref{c:rungs.transported-module-holds-tuple-one-vector}); and the dimension counts of Theorem \ref{c:rungs.rung-spinor-module-mixings}(4). What the script does not check: the Facts beyond the instances named, verified in their own papers' scripts; the real-form statement of Remark \ref{c:note.rungs.compact-form}, which is representation theory imported through Fact \ref{c:rungs.imported-fock-representation-bogoliubov-group}; and Hypothesis (R), which is not a computation.

\section{Results used from elsewhere}
The proofs above use the following results of the series, stated briefly; each is proved in the paper named.

\begin{fact}[The Fock model of the complexified ascended algebra]\label{c:correspondence.fock-model-complexified-ascended-algebra}
The exterior algebra on a maximal isotropic subspace, with wedge and contraction operators, is the unique irreducible module of the complexified Clifford algebra $\Cl(2,4)_\C$, and its even and odd halves are the two four-dimensional chiralities.
\emph{Source:} T9, Thm 3.1 (\href{https://doi.org/10.5281/zenodo.22722290}{doi:10.5281/zenodo.22722290}).
\end{fact}

\begin{fact}[The incidence relation is the chart]\label{c:correspondence.incidence-relation-chart}
The Pl\"ucker bivector of the incidence line of a map $A$ is exactly the completion's null point embedding $X(m(A))$, so Penrose's incidence relation is the completion's affine chart, with a determinant as its quadratic term.
\emph{Source:} T9, Thm 4.2 (\href{https://doi.org/10.5281/zenodo.22722290}{doi:10.5281/zenodo.22722290}).
\end{fact}

\begin{fact}[Reality: the conjugation, the invariant Hermitian form, and the factor $I_6$]\label{c:correspondence.reality-conjugation-twistor-form-factor}
The completion carries an antilinear conjugation, unique up to scale, whose invariant form on the conformal spinor space is Hermitian of signature $(2,2)$, and the real chart points are exactly the null lines of that form skew to the infinity line, with incidence relation $\omega = I_6\hat x\pi$ for Hermitian $\hat x$.
\emph{Source:} T9, Prop 4.5 (\href{https://doi.org/10.5281/zenodo.22722290}{doi:10.5281/zenodo.22722290}).
\end{fact}

\begin{fact}[Skew lines, splittings, hyperbolic planes: one datum]\label{c:correspondence.skew-lines-splittings-hyperbolic-planes-one}
An ordered pair of skew lines in the projective spinor space, a normalised splitting of the conformal spinor space into two planes and a normalised null pair in $V_\C$ determine one another, with the mixed slot equal to the orthogonal complement of the null pair.
\emph{Source:} T9, Prop 4.1 (\href{https://doi.org/10.5281/zenodo.22722290}{doi:10.5281/zenodo.22722290}).
\end{fact}

\begin{fact}[The spinor identification]\label{c:correspondence.spinor-identification}
There is a canonical $\Spin$-equivariant isomorphism from $V_\C$ to the bivectors of the conformal spinor space carrying the quadratic form to the wedge square, so null vectors are exactly the simple bivectors, that is, the lines of the projective spinor space.
\emph{Source:} T9, Thm 3.4 (\href{https://doi.org/10.5281/zenodo.22722290}{doi:10.5281/zenodo.22722290}).
\end{fact}

\begin{fact}[The de Sitter scale and its frozen point]\label{c:mass.de-sitter-scale-frozen-point}
On flat spacetime de Sitter space is the gravitational field $\hb = \sigma$, $\sigma = 1 - \Lambda Q(x)/12$, with derived metric $\sigma^{-2}\eta$, and its frozen point is the unique constant vector $I_\Lambda = n_\infty + (\Lambda/6)n_0$, of squared norm $\Lambda/3$, which is the point at infinity $n_\infty$ at $\Lambda = 0$.
\emph{Source:} T11, Lem 5.1 (\href{https://doi.org/10.5281/zenodo.22722565}{doi:10.5281/zenodo.22722565}).
\end{fact}

\begin{fact}[Mass as symplectic area: the two-spinor formula]\label{c:mass.mass-symplectic-area-two-twistor-formula}
Two conformal spinors with momentum $\pi^{1}\bar\pi^{1} + \pi^{2}\bar\pi^{2}$ have $m^{2} = 2|\br{\pi^{1}\pi^{2}}|^{2}$, the squared symplectic area of their spinor parts.
\emph{Source:} T11, Prop 4.1 (\href{https://doi.org/10.5281/zenodo.22722565}{doi:10.5281/zenodo.22722565}).
\end{fact}

\begin{fact}[The gauge algebra of the minimal massive composite]\label{c:standard-model.gauge-algebra-minimal-massive-composite}
The automorphism algebra of the minimal massive hierarchy $\C^{2} \otimes \C^{3}$ is $\su(2) \oplus \su(3) \oplus \uu(1)$, exactly the Standard Model's, while that of $\C^{3} \otimes \C^{3}$ is seventeen-dimensional.
\emph{Source:} T12, Thm 4.1 (\href{https://doi.org/10.5281/zenodo.22722669}{doi:10.5281/zenodo.22722669}).
\end{fact}

\begin{fact}[The global form: the quotient determined]\label{c:standard-model.global-form-quotient-determined}
The global gauge group of the minimal massive composite is $S(U(2) \times U(3)) \cong (SU(2) \times SU(3) \times U(1))/\mathbb{Z}_{6}$, the smallest of the Standard Model's global forms.
\emph{Source:} T12, Thm 4.2 (\href{https://doi.org/10.5281/zenodo.22722669}{doi:10.5281/zenodo.22722669}).
\end{fact}

\begin{fact}[Hypothesis (S): slot statistics, derived from the slot-module hypothesis]\label{c:standard-model.hyp-s-slot-statistics}
Under the slot-module hypothesis, the composite's slots are fermionic modes obeying exclusion and forming one Clifford system with its pairs, its state space is the five-mode Fock space $\Lambda^{\bullet}(\C^{2} \oplus W)$ with $U(2) \times U(3)$ acting as $\Lambda^{\bullet}(A \oplus B)$, chirality is the parity of the mode number, and slots of one tower parity whose grading is the chirality are odd.
\emph{Source:} T12, Hypothesis (S) (\href{https://doi.org/10.5281/zenodo.22722669}{doi:10.5281/zenodo.22722669}).
\end{fact}

\begin{fact}[Hypercharge from the global form, pair counting, and consistency]\label{c:standard-model.hypercharge-global-form-pair-counting-consistency}
Under Hypothesis (S) for the doublets, the composite's states carry the observed hypercharges of one generation, which satisfy all four anomaly conditions and are the unique solution for that content with $Y_{Q} = \tfrac16$.
\emph{Source:} T12, Thm 5.5 (\href{https://doi.org/10.5281/zenodo.22722669}{doi:10.5281/zenodo.22722669}).
\end{fact}

\begin{fact}[Imported: the Fock representation of a Clifford algebra, and its Bogoliubov group]\label{c:standard-model.imported-fock-representation-clifford-algebra-bogoliubov}
The Fock space $\Lambda^{\bullet}E$ of $n$ fermionic modes is the irreducible spinor module of the complex Clifford algebra on $2n$ generators, with $\mathrm{Pin}(2n)$ implementing the Bogoliubov group and $U(n)$ the number-preserving part.
\emph{Source:} W. Fulton and J. Harris, Representation Theory: A First Course, Graduate Texts in Mathematics 129, Springer (1991)., Lecture~20, \S 20.1; H. B. Lawson and M.-L. Michelsohn, Spin Geometry, Princeton University Press (1989)., Ch.~I, \S\S 2, 5.
\end{fact}

\begin{fact}[Imported: the frozen point and the dictionary]\label{c:standard-model.imported-frozen-point-dictionary}
Mass is conformal weight read against the frozen point, $m^{2} = \Delta(3-\Delta)H^{2}$, so heavy fields have complex weight and the Planck mass's imaginary weight is $N \approx M_{P}/H$, about $10^{61}$.
\emph{Source:} T12, Fact 2.5 (\href{https://doi.org/10.5281/zenodo.22722669}{doi:10.5281/zenodo.22722669}).
\end{fact}

\begin{fact}[The internal symmetry ladder]\label{c:standard-model.internal-symmetry-ladder}
The mixings of a system of $n$ conformal spinors that preserve total momentum and total helicity form exactly $U(n)$, giving the ladder $U(1)$, $U(2)$, $U(3)$.
\emph{Source:} T12, Prop 3.1 (\href{https://doi.org/10.5281/zenodo.22722669}{doi:10.5281/zenodo.22722669}).
\end{fact}

\begin{fact}[The Majorana element is the $(10,1,3)$]\label{c:standard-model.majorana-element}
Under Hypothesis (S), the product of the five creation operators spans the $(10,1,3)$ under Pati-Salam, gives only $\nu^{c}$ a Majorana mass, and is fixed by exactly the Standard Model algebra.
\emph{Source:} T12, Prop 10.7 (\href{https://doi.org/10.5281/zenodo.22722669}{doi:10.5281/zenodo.22722669}).
\end{fact}

\begin{fact}[The mass bracket is a charged singlet]\label{c:standard-model.mass-bracket-charged-singlet}
The mass bracket transforms by $\det(U)$ under the two-spinor $U(2)$, making it an $SU(2)$ singlet of hypercharge $+1$ with the quantum numbers of $e^{c}$, not of a weak doublet.
\emph{Source:} T12, Prop 3.2 (\href{https://doi.org/10.5281/zenodo.22722669}{doi:10.5281/zenodo.22722669}).
\end{fact}

\begin{fact}[Quarks and leptons: the two-level fermion content]\label{c:standard-model.quarks-leptons-two-level-fermion-content}
The pair-level cells carry the quark doublet $(2,3)_{1/6}$ unconditionally, and under Hypothesis (S) one slot with the pairs full gives the lepton doublet $(2,1)_{-1/2}$.
\emph{Source:} T12, Prop 5.3 (\href{https://doi.org/10.5281/zenodo.22722669}{doi:10.5281/zenodo.22722669}).
\end{fact}

\begin{fact}[The spectrum run]\label{c:standard-model.spectrum-run}
At one loop with the record inputs, the derived content with one imported doublet predicts $\sin^{2}\theta_{W}(M_{Z}) = 78761/379488 \approx 0.2075$ against $0.2312$, refuting the desert, while the supersymmetric coefficients give $0.2309$; the framework's $3/8$ survives only if the running differs from that of the derived content.
\emph{Source:} T12, Prop 6.2 (\href{https://doi.org/10.5281/zenodo.22722669}{doi:10.5281/zenodo.22722669}).
\end{fact}

\begin{fact}[Under (S): the composite's state space is the $\mathrm{Spin}(10)$ spinor, and the derived generation is its even half]\label{c:standard-model.under-s-composites-state-space-spinor}
Under Hypothesis (S), the composite's state space $\Lambda^{\bullet}(\C^{2} \oplus \C^{3})$ is the $\mathrm{Spin}(10)$ spinor, its even half is the derived generation, and the extra generator $U(1)_{X}$ is anomaly-free in the action it carries.
\emph{Source:} T12, Prop 5.8 (\href{https://doi.org/10.5281/zenodo.22722669}{doi:10.5281/zenodo.22722669}).
\end{fact}

\begin{fact}[The weak angle over the derived generation]\label{c:standard-model.weak-angle-over-derived-generation}
Over one derived generation, with the lepton under Hypothesis (S), $\sum T_{3}^{2}/\sum Q^{2} = 3/8$, which is $\sin^{2}\theta_{W}$ at unification under the simple-group hypothesis.
\emph{Source:} T12, Prop 6.1 (\href{https://doi.org/10.5281/zenodo.22722669}{doi:10.5281/zenodo.22722669}).
\end{fact}

\begin{fact}[Complexification, and the odd-dimensional split]\label{c:std-algebra.complexification}
$\Cl(V) \otimes \C \cong \Cl_n(\C)$, a complex matrix algebra (doubled for odd $n$); for odd $n$, $\Cl \cong \Cl^{+} \oplus \Cl^{+}$ if $\omega^{2} = 1$ and $\Cl \cong \Cl^{+} \otimes \C$ if $\omega^{2} = -1$.
\emph{Source:} the algebra core of the Transport Lab, proved there with an exact check; standard, see P. Lounesto, Clifford Algebras and Spinors, 2nd ed., London Mathematical Society Lecture Note Series 286, Cambridge University Press (2001), Ch. 16 (matrix representations and periodicity of 8).
\end{fact}

\begin{fact}[Commuting copies]\label{c:tower.commuting-copies}
Each floor $\Cl(p{+}k, q{+}k)$ of the tower is canonically a twisted copy of $\Cl(p,q)$ tensored with $\Mat_{2^k}(\R)$, that is $\Mat_{2^k}(\Cl(p,q))$.
\emph{Source:} T3, Thm 3.2 (\href{https://doi.org/10.5281/zenodo.22721556}{doi:10.5281/zenodo.22721556}).
\end{fact}

\begin{fact}[Fock module; one qubit per rung]\label{c:tower.fock-module-one-qubit-per-rung}
The Chevalley action makes the exterior algebra $\Lambda^\bullet W_k$ an irreducible $2^k$-dimensional module for $B_k \cong \Mat_{2^k}(\R)$ with a cyclic vacuum, and each rung doubles it.
\emph{Source:} T3, Thm 4.2 (\href{https://doi.org/10.5281/zenodo.22721556}{doi:10.5281/zenodo.22721556}).
\end{fact}

\begin{fact}[The ladder pairs satisfy the CAR, as vectors]\label{c:tower.ladder-pairs-satisfy-car-vectors}
In the $k$-th tower algebra the $k$ adjoined null pairs satisfy the canonical anticommutation relations of $k$ fermionic modes exactly, as vectors, and each new rung adds one mode without altering the existing relations.
\emph{Source:} T3, Prop 4.1 (\href{https://doi.org/10.5281/zenodo.22721556}{doi:10.5281/zenodo.22721556}).
\end{fact}

\begin{fact}[The *-structure of the tower, and its C*-norms]\label{c:tower.structure-tower-c-norms}
Each complexified floor of the tower has a unique frame-adapted *-structure respected by all the tower's maps and correctors, and for $p - q$ even each floor is a full complex matrix algebra, so the floors and their union carry unique C*-norms.
\emph{Source:} T3, Lem 5.4 (\href{https://doi.org/10.5281/zenodo.22721556}{doi:10.5281/zenodo.22721556}).
\end{fact}

\begin{fact}[Witt index is quantum capacity]\label{c:tower.witt-index-quantum-capacity}
For a nondegenerate real quadratic space of signature $(p,q)$, the maximal number of fermionic mode pairs realisable by vectors of its Clifford algebra is the Witt index $\min(p,q)$, which rises by exactly one at each rung of the tower.
\emph{Source:} T3, Thm 4.3 (\href{https://doi.org/10.5281/zenodo.22721556}{doi:10.5281/zenodo.22721556}).
\end{fact}

\begin{fact}[Minimality]\label{c:transport.minimality}
Every null-lift completion adjoins directions of both signs, so the minimal completion is the hyperbolic-plane ascent $\Cl(p,q) \to \Cl(p{+}1,q{+}1)$, unique up to isometry, with increment $(+1,+1)$.
\emph{Source:} T1, Thm 4.13 (\href{https://doi.org/10.5281/zenodo.21743863}{doi:10.5281/zenodo.21743863}).
\end{fact}

\begin{fact}[No-go at the level of elements]\label{c:transport.no-go-level-elements}
In mixed signature with $n \geq 3$, no assignment of von Neumann pseudo-inverses to null vectors, valued in the Clifford algebra of any extension, is equivariant even under $\Spin^0(p,q)$.
\emph{Source:} T1, Thm 4.2 (\href{https://doi.org/10.5281/zenodo.21743863}{doi:10.5281/zenodo.21743863}).
\end{fact}

\begin{fact}[Totality]\label{c:transport.totality}
On the quadric of the completion, $\iota[N] = [e_+ N e_+]$ is a total continuous involution that equals the reciprocal $x \mapsto x/Q(x)$ on invertible vectors and sends each $x_0$ with $Q(x_0) = 0$ to $[x_0 + \frac12 \nInf]$, and it is the only such continuous extension.
\emph{Source:} T1, Thm 4.10 (\href{https://doi.org/10.5281/zenodo.21743863}{doi:10.5281/zenodo.21743863}).
\end{fact}

\begin{refutation}[The desert: the derived content alone connects the structural scale to experiment]\label{c:refuted.the-desert-the-derived-content-alone-connects-th}
Refuted: The desert: the derived content alone connects the structural scale to experiment
\end{refutation}

\begin{refutation}[The internal ladder acts on the whole tower at once]\label{c:refuted.the-internal-ladder-acts-on-the-whole-tower-at-o}
Refuted: The internal ladder acts on the whole tower at once
\end{refutation}

\begin{refutation}[The raw representation content matches the fermion table entry by entry]\label{c:refuted.the-raw-representation-content-matches-the-fermi}
Refuted: the raw representation content does not match the fermion table entry by entry; the disagreements are composites, the generators are the content, and the lepton lies one level higher under the slot-statistics hypothesis.
\end{refutation}

\begin{refutation}[The hierarchy of three conformal spinors $\mathbb{C}^{3}\otimes\mathbb{C}^{3}$ carries the Standard Model]\label{c:refuted.the-three-twistor-hierarchy-carries-the-standard}
Refuted: the hierarchy of three conformal spinors $\C^{3}\otimes\C^{3}$ has a seventeen-dimensional automorphism algebra, not the Standard Model's twelve, which the minimal massive composite $\C^{2}\otimes\C^{3}$ carries instead.
\end{refutation}

\begin{refutation}[The original conjecture, refuted below]\label{c:transport.original-conjecture-refuted-below}
Element-by-element repair of a null vector by an equivariant partner one ascent up: refuted.
\end{refutation}

\begin{prediction}[No fractionally charged colour singlet]\label{c:pred.z6-charge-lattice}
The $\mathbb{Z}_{6}$ charge lattice of the derived gauge group forbids any fractionally charged colour singlet.
\emph{Kill condition:} One confirmed fractionally charged colour-singlet particle.
\end{prediction}

\begin{thebibliography}{99}
\bibitem{ButlerT12} L.~Butler,
\emph{The Standard Model from the Minimal Massive Composite},
Transport Papers T12 (2026), doi:\href{https://doi.org/10.5281/zenodo.22722669}{10.5281/zenodo.22722669}.
\bibitem{ButlerT3} L.~Butler,
\emph{The Tower and the CAR Limit: Iterated Transport, the Frozen Clock, and Second Quantisation},
Transport Papers T3 (2026), doi:\href{https://doi.org/10.5281/zenodo.22721556}{10.5281/zenodo.22721556}.
\bibitem{ButlerT1} L.~Butler,
\emph{The Transport Principle: Total Division in Clifford Algebras via Minimal Completion},
Transport Papers T1 (2026).
\bibitem{ButlerT9} L.~Butler,
\emph{The Twistor Correspondence: Ascended Spinors, Incidence, and the Palatial Bridge},
Transport Papers T9 (2026).
\bibitem{ButlerT11} L.~Butler,
\emph{The Origin of Mass: The Weight Dictionary and the Exact Join},
Transport Papers T11 (2026), doi:\href{https://doi.org/10.5281/zenodo.22722565}{10.5281/zenodo.22722565}.
\bibitem{FultonHarris1991} W.~Fulton and J.~Harris,
\emph{Representation Theory: A First Course},
Graduate Texts in Mathematics 129, Springer (1991).
\bibitem{LawsonMichelsohn1989} H.~B. Lawson and M.-L. Michelsohn,
\emph{Spin Geometry},
Princeton University Press (1989).
\bibitem{ButlerT17} L.~Butler,
\emph{Capstone: The Transport Programme, Method, Status, and the Falsifiable Ledger},
Transport Papers T17 (2026).
\bibitem{ButlerVerifyT12R} L.~Butler,
\emph{verify\_T12\_rungs.py: machine verification for ``Two Rungs Up''},
Python script, deposited with this note alongside the shared library GA.py (2026).
\end{thebibliography}
\end{document}